\begin{abstract}
Digital government services frequently provide citizen functionalities across multiple channels, such as web portals and mobile applications, which may exhibit asymmetries in data completeness, response time, and operational resilience. This paper presents the definition and specification phases (Phases 1 and 2 of ISO/IEC 25040) for a comparative black-box software quality evaluation of the traffic infraction consultation service provided by Detran-DF across its Services Portal and mobile application (Android/iOS). By adapting the ISO/IEC 25010:2023 and ISO/IEC 25010:2011 quality models, we prioritize functional suitability, performance efficiency, and context coverage. Using the Goal Question Metric (GQM) approach, we derive three measurement goals, seven analytical questions, falsifiable hypotheses, and fourteen external metrics with objective scoring levels and judgment criteria, alongside project governance tracked via Practical Software Measurement (PSM/CID).
\end{abstract}

\begin{resumo}
Serviços digitais governamentais frequentemente disponibilizam funcionalidades ao cidadão em múltiplos canais, como portais web e aplicativos móveis, os quais podem apresentar assimetrias de completude informacional, tempo de resposta e resiliência operacional. Este artigo apresenta a definição dos requisitos e a especificação da avaliação (Fases 1 e 2 da norma ISO/IEC 25040) para uma avaliação comparativa em caixa-preta da qualidade do serviço de consulta de infrações de trânsito do Detran-DF entre o Portal de Serviços e o aplicativo móvel oficial (Android/iOS). A partir da adaptação dos modelos de qualidade da ISO/IEC 25010:2023 e da ISO/IEC 25010:2011, priorizam-se as características de adequação funcional, eficiência de desempenho e cobertura de contexto. Empregando a abordagem Goal Question Metric (GQM), derivam-se três objetivos de medição, sete questões analíticas, hipóteses verificáveis e quatorze métricas externas com níveis de pontuação e critérios objetivos de julgamento, apoiados por governança baseada no Practical Software Measurement (PSM/CID).
\end{resumo}